%======================================================================
\section{Problem setting}
%======================================================================

Users $u\in[m]$, items $i\in[n]$ and pairs $e=(u,i)$ carry a binary potential
outcome $Y_e\in\{0,1\}$ (whether $u$ would like $i$ if it were shown) and an
observation indicator $O_e\in\{0,1\}$. The log reveals $O$ and $O\odot Y$. What
$O$ records depends on the platform: on KuaiRec it is an impression chosen by
the system, on Coat and Yahoo!\,R3 it is a rating that the user chose to give,
so the ``exposure'' mechanism modelled by $p$ is then user self-selection
rather than the platform's display policy. All statements below are relative
to the mechanism that generates $O$; we keep the word exposure for brevity.

\begin{assumption}[Conditionally independent exposure]\label{as:unconf}
Conditionally on $Y$, the exposures $O_e$ are independent Bernoulli variables
with propensities $p_e=\Pr(O_e=1\mid Y)>0$.
\end{assumption}
The propensity may depend on the pair's own outcome, so the log may be missing
not at random. All expectations below are over $O$ given $Y$.

\begin{assumption}[Fixed nuisances]\label{as:fixed}
The propensity estimates $\hat p$, the imputation $\hat Y\in[0,1]^{m\times n}$,
the edge weights $C$ and the constants $a,b$ of the score below are fixed
given $Y$, i.e., they do not depend on $O$ (for instance, they are fitted on
an independent log). We use $C_{ui}=d_u^{-\alpha}d_i^{-(1-\alpha)}$ with degrees
$d$ taken from $\hat Y$.
\end{assumption}
The weights $C$ also enter the target below, so with degrees from $\hat Y$ the
target is indexed by the imputation. A normalisation fixed a priori (a design
choice, or degrees from an independent source) makes the target independent
of $\hat Y$; Section~\ref{sec:semisynth} compares both.

\begin{remark}[Scope of the guarantees and the experimental protocol]\label{rem:xfit}
All guarantees below are stated under Assumption~\ref{as:fixed}. Our main
experimental protocol does not satisfy it, in four ways. (i) $\hat p$ and
$\hat Y$ are cross-fitted over ten random folds of pairs: the nuisances used at
a pair of fold $k$ are fitted on the other nine
folds~\cite{chernozhukov2018dml}. This makes the nuisance at a pair
independent of that pair's exposure, which is what a one-hop estimate needs,
but not of the exposures of the other edges of a walk, so a product of
several edge estimates can be biased although every single edge estimate is
unbiased. (ii) The degree weights sum the imputation over all pairs of a user
or an item, so they depend on the exposures of all these pairs, including the
pair at which they are used. (iii) One of the searched degree sources takes
degrees from edge estimates. (iv) The constants $a,b$ are computed from the
scores (Section~\ref{sec:drup}). In simulation, the resulting dependence bias
is of the same order as the repeated-walk bias, and under this protocol with
imputation-based degrees the corrected operator is not less biased than the
uncorrected one
(Section~\ref{sec:mc}); a control in which only the degree weights are fixed a
priori restores the correction under cross-fitting, so (ii) and (iii), more than (i),
account for the failure. The theorems therefore describe DRUP with nuisances
from an independent log, or from the sample split introduced below; for the
cross-fitted protocol they give no guarantee, and our experiments with it
evaluate accuracy and the empirical behaviour of the scores, not unbiasedness.
\end{remark}

With $\tY=C\odot Y$, the target is propagation over the full-exposure graph,
\begin{equation}\label{eq:target}
F^*=a\,\tY+b\,T^*,\qquad T^*=\tY\tY^\top\tY ,
\end{equation}
with fixed constants $a,b\ge0$. $F^*$ is the linear LightGCN score one would
compute if every item had been shown to every user. It is a normative choice:
it treats potential outcomes as fixed and ignores preferences that exposure
itself shapes~\cite{chaney2018confounding,perdomo2020performative}
(Section~\ref{sec:disc}).

%======================================================================
\section{DRUP}\label{sec:drup}
%======================================================================

\paragraph{Edge estimates.}
With clipped propensities $\bar p=\max(\hat p,\tau)$,
\begin{equation}
W^{\mathrm{IPS}}_e=\frac{O_eY_e}{\bar p_e},\qquad
W^{\mathrm{DR}}_e=\hat Y_e+\frac{O_e(Y_e-\hat Y_e)}{\bar p_e}.
\end{equation}
Both have mean $Y_e$ when $\bar p=p$.

\paragraph{Where the plug-in operator fails.}
The plug-in three-hop operator $(\tW\tW^\top\tW)_{ui}$, with $\tW=C\odot W$, sums
the products $\tW_{uj}\tW_{vj}\tW_{vi}$ over all walks $u\!-\!j\!-\!v\!-\!i$.
If $v\neq u$ and $j\neq i$ the three edges are distinct and, by independence,
the product is unbiased. The other walks repeat an edge
(Figure~\ref{fig:walks}): for $v=u$ the edge $(u,j)$ is traversed twice, for
$j=i$ the edge $(v,i)$, and for $v=u,j=i$ the edge $(u,i)$ three times. A
repeated factor $\tW_e^k$ has mean $C_e^k\E[W_e^k]\ne C_e^kY_e$; for IPS,
$\E[W_e^2]=Y_e/p_e$. Because $Y_e^k=Y_e$ for binary outcomes, the unbiased
replacement of $\tW_e^k$ is simply $C_e^kW_e$.

\begin{figure}[t]\centering
\begin{tikzpicture}[x=1cm,y=1cm,font=\footnotesize,
  un/.style={circle,draw,minimum size=5.5mm,inner sep=0pt,fill=blue!8},
  it/.style={rectangle,draw,minimum size=5mm,inner sep=0pt,fill=orange!12},
  w/.style={->,thick,>=stealth},
  rep/.style={->,very thick,>=stealth,red!75!black}]
\begin{scope}[xshift=0cm]
 \node[un] (u) at (0,1) {$u$}; \node[un] (v) at (0,0) {$v$};
 \node[it] (j) at (1.6,1) {$j$}; \node[it] (i) at (1.6,0) {$i$};
 \draw[w] (u) -- (j); \draw[w] (j) -- (v); \draw[w] (v) -- (i);
 \node at (0.8,-0.75) {(a) distinct edges};
\end{scope}
\begin{scope}[xshift=3.3cm]
 \node[un] (u) at (0,1) {$u$};
 \node[it] (j) at (1.6,1) {$j$}; \node[it] (i) at (1.6,0) {$i$};
 \draw[rep] (u) to[bend left=15] node[above]{$\times2$} (j); \draw[rep] (j) to[bend left=15] (u); \draw[w] (u) -- (i);
 \node at (0.8,-0.75) {(b) $v=u$: $\tW_{uj}^2\tW_{ui}$};
\end{scope}
\begin{scope}[xshift=6.6cm]
 \node[un] (u) at (0,1) {$u$}; \node[un] (v) at (0,0) {$v$};
 \node[it] (i) at (1.6,0.5) {$i$};
 \draw[w] (u) -- (i); \draw[rep] (i) to[bend left=15] node[below right=-1pt]{$\times2$} (v); \draw[rep] (v) to[bend left=15] (i);
 \node at (0.8,-0.75) {(c) $j=i$: $\tW_{ui}\tW_{vi}^2$};
\end{scope}
\begin{scope}[xshift=9.9cm]
 \node[un] (u) at (0,0.5) {$u$};
 \node[it] (i) at (1.6,0.5) {$i$};
 \draw[rep] (u) to[bend left=25] node[above]{$\times3$} (i); \draw[rep] (i) -- (u); \draw[rep] (u) to[bend right=25] (i);
 \node at (0.8,-0.75) {(d) $v=u$, $j=i$: $\tW_{ui}^3$};
\end{scope}
\end{tikzpicture}
\caption{The four kinds of three-hop walks from user $u$ to item $i$. Only walks
of type (a) consist of distinct edges. In (b)--(d) an edge estimate is
multiplied with itself (red), and its power has the wrong mean. DRUP replaces
each repeated factor $\tW_e^k$ by $C_e^kW_e$.}\label{fig:walks}
\end{figure}

\paragraph{Idempotent walk correction.}
Let $\Delta=\tW\odot\tW-C\odot C\odot W$, with row sums $r$ and column sums
$\kappa$. The corrected three-hop operator subtracts the excess of each walk
type:
\begin{equation}\label{eq:drup}
\hat T=\tW\tW^\top\tW-\tW\odot(r\mathbf 1^\top-\Delta)-\tW\odot(\mathbf 1\kappa^\top-\Delta)
-(\tW^{\odot3}-C^{\odot3}\odot W).
\end{equation}
The second term removes the excess of type (b) walks, the third that of type
(c), and the last that of type (d); the subtracted $\Delta$ inside the first
two avoids counting type (d) twice. The DRUP score is $s=a\,\tW+b\,\hat T$.
In the experiments $a=1/c_1$ and $b=\beta/c_3$, where $c_1$ and $c_3$ are the
mean absolute values of the one- and three-hop terms over all scored users and
$\beta$ is selected on validation data. The constants are global, so every
user's scores are the same linear combination of the two terms, and
within-user rankings depend on them only through $b/a$; a data-dependent
$c_3/c_1$ thus acts as a data-dependent rescaling of the grid of $\beta$.
Because $c_1$ and $c_3$ are computed from the log, however,
$\E[\hat T/c_3]\neq T^*/c_3$ in general, and the score-level statements below
hold for fixed $a,b$. Constants computed from the imputed graph $C\odot\hat Y$
satisfy Assumption~\ref{as:fixed} whenever $\hat Y$ does; Section~\ref{sec:mc}
compares both choices.

\paragraph{Sample-split variant.}
A log rarely comes with an independent log for its nuisances. A sample split
provides one at a cost. Draw a random set $A$ of pairs, each pair with
probability $q$, fit $\hat Y$ (and $\hat p$, where it is estimated) on the
pairs in $A$ only, replace
the edge estimates of the pairs in $A$ by $\hat Y$, take the degrees from
$\hat Y$, and compute $c_1$ and $c_3$ on the imputed graph $C\odot\hat Y$.
Conditionally on $A$ and on the exposures in $A$, the remaining edge estimates
are independent of every nuisance, so Assumption~\ref{as:fixed} holds and the
results below apply, with $\xi_e=\hat Y_e-Y_e$ on the pairs in $A$. The price
is a first-order imputation bias on a fraction $q$ of the pairs and fewer
exposures for the edge estimates. We call this variant DRUP-split and use
$q=0.2$. The guarantee holds for every fixed configuration; selecting the
hyper-parameters on validation data, as every method in our comparison does,
is outside it, as it is for any estimator.

\paragraph{Public-operator form.}
Define the diagonal-corrected item Gram matrix
$G=\tW^\top\tW-\diag(\tW^\top\tW)+\diag(\mathbf 1^\top(C\odot C\odot W))$. Then
\begin{equation}\label{eq:local}
\hat T_{u\cdot}=\tilde w_uG-\tilde w_u\odot(r_u\mathbf 1^\top-2\zeta_u)-(\tilde w_u^{\odot3}-c_u^{\odot3}\odot w_u),
\end{equation}
where $\zeta_u=\tilde w_u^{\odot2}-c_u^{\odot2}\odot w_u$ and lower-case rows
belong to user $u$. Equation~\eqref{eq:local} matches~\eqref{eq:drup} to machine
precision. $G$ is the only object shared across users
(Algorithm~\ref{alg:drup}). Its off-diagonal entries are products of distinct
edges; its diagonal is corrected from $\sum_u\tW_{ui}^2$ to
$\sum_uC_{ui}^2W_{ui}$, the analogue of the diagonal correction of a Gram
matrix estimated from masked data~\cite{lounici2014}.

\begin{algorithm}[t]
\caption{DRUP (training-free)}\label{alg:drup}
\begin{algorithmic}[1]
\Require log $(O,O\odot Y)$; hyper-parameters $\tau,\alpha,\lambda,\beta$
\State fit $\hat p$ (logistic exposure model) and $\hat Y$ (IPS-weighted additive ridge, optionally plus an IPS-weighted low-rank ALS residual), cross-fitted over ten folds of pairs
\State $\bar p\gets\max(\hat p,\tau)$; \ $W\gets\hat Y+O\odot(Y-\hat Y)/\bar p$
\State $C\gets d_u^{-\alpha}d_i^{-(1-\alpha)}$ with $d$ from $\hat Y$
\State $G\gets\tW^\top\tW-\diag(\tW^\top\tW)+\diag(\mathbf 1^\top(C\odot C\odot W))$ \Comment{public operator}
\For{each user $u$ (locally)}
\State $\hat T_{u\cdot}\gets$ Eq.~\eqref{eq:local}; \ $s_u\gets a\,\tilde w_u+b\,\hat T_{u\cdot}$; rank unexposed items by $s_u$
\EndFor
\end{algorithmic}
\end{algorithm}

\paragraph{Control-variate family.}
More generally, $W^\lambda_e=O_eY_e/\bar p_e+\lambda(\hat Y_e-O_e\hat Y_e/\bar p_e)$
is unbiased for every $\lambda\in[0,1]$ under correct propensities. It
interpolates between IPS ($\lambda=0$) and DR ($\lambda=1$, the only member that
is also robust to propensity errors), and it equals the DR estimate with the
shrunk imputation $\lambda\hat Y$. With $\bar p=p$ its variance is
$\frac{1-p}{p}(Y-\lambda\hat Y)^2$. All results below hold for every member, and
$\lambda$ is selected on validation data.

\paragraph{Imputation.}
Theorem~\ref{thm:unbiased} holds for any fixed $\hat Y$, so the imputation is a
design choice; its accuracy enters the variance and the clipping bias
(Theorem~\ref{thm:var}). We use an additive model $\mu+b_u+b_i$ fitted by
IPS-weighted ridge regression, or that model plus a rank-32 residual $UV^\top$
fitted by IPS-weighted alternating least squares. The additive imputation gives
every user the same item order on unexposed candidates, which makes the one-hop
DR term non-personalised; the low-rank imputation removes this limitation.

\paragraph{Complexity and scale.}
At three hops, one $n\times n$ Gram product and $O(mn)$ reductions, i.e., the
cost of the uncorrected operator. The $K$-hop correction of
Theorem~\ref{thm:khop} evaluates one tensor contraction per pair of equality
patterns; we report its measured cost for $K=5$ in Section~\ref{sec:mc}. On KuaiRec (7{,}176 users, 10{,}728 items) one
configuration takes a few seconds on four CPU cores once the nuisances are
fitted. The dense $W$ and $G$
need $O(mn+n^2)$ memory, which limits the present implementation to catalogues
of about $10^4$--$10^5$ items. Because the correction only modifies diagonals
and row and column sums, it carries over to sparse or low-rank
representations of $W$, but we have not evaluated catalogues of $10^6$ items.

%======================================================================
\section{Theory}
%======================================================================

Proofs are in Appendix~\ref{app:proofs}. Let $\varepsilon_e=|Y_e-\hat Y_e|$.

\subsection{Unbiasedness and its limits}

\begin{theorem}[Edge-wise double robustness]\label{thm:unbiased}
Under Assumptions~\ref{as:unconf}--\ref{as:fixed}, let
$\xi_e=(p_e/\bar p_e-1)(Y_e-\hat Y_e)$. Then $\E[\hat T_{ui}]$ equals $T^*_{ui}$
with every distinct edge contributing $Y_e+\xi_e$ once. In particular, if
every edge has a correct propensity or a correct imputation, then
$\E[s]=F^*$.
\end{theorem}

Recommendations rank the \emph{unexposed} candidates of a user, and there
the candidate's own edge estimate is deterministic ($W_{ui}=\hat Y_{ui}$ for DR,
$0$ for IPS). The following corollary states what DRUP estimates on the
candidate set.

\begin{corollary}[Unexposed candidates]\label{cor:cand}
Under the conditions of Theorem~\ref{thm:unbiased} with $\xi\equiv0$,
\begin{gather*}
\E[s_{ui}\mid O_{ui}=0]=F^*_{ui}-(Y_{ui}-\hat Y_{ui})\,g_{ui},\\
g_{ui}=\frac{\partial F^*_{ui}}{\partial Y_{ui}}
=C_{ui}\Big(a+b\big(\textstyle\sum_{v\neq u}C_{vi}^2Y_{vi}+\sum_{j\neq i}C_{uj}^2Y_{uj}+C_{ui}^2\big)\Big)\ge0 .
\end{gather*}
For IPS edges the same holds with $\hat Y_{ui}=0$. The residual is the
imputation error at the candidate itself. If $p_e$ depends on $Y$ only through
$Y_e$, the log carries no information on $Y_{ui}$ when $O_{ui}=0$, so the
residual is not identified and no estimator can remove it.
\end{corollary}

At three hops the walks through $(u,i)$ are exactly those counted in
$g_{ui}$. With IPS edges they vanish on unexposed candidates, so the uncorrected
IPS operator has no repeated-walk bias there; from five hops on, walks repeat
edges that are not incident to the candidate (for example
$u\!-\!j\!-\!v\!-\!j\!-\!v\!-\!i$) and the bias reappears
(Table~\ref{tab:mc}).

\begin{theorem}[Any odd number of hops]\label{thm:khop}
For $K=2r+1$, a walk $U_0\!-\!J_1\!-\!U_1\!-\cdots-\!U_r\!-\!J_{r+1}$ has user
variables $U_0,\dots,U_r$ and item variables $J_1,\dots,J_{r+1}$, with $U_0=u$
and $J_{r+1}=i$. The equality pattern of an assignment is a pair $\pi$ of set
partitions of the user and item variables, and $\pi$ fixes the multiplicity
$k_e$ of every distinct edge. Let $g_\pi(\sigma)$ be the unconstrained
contraction, over the blocks of a coarser pattern $\sigma\ge\pi$, of
$\prod_eC_e^{k_e}W_e$ (corrected) or of $\prod_e\tW_e^{k_e}$ (plug-in), the
product running over the distinct edges of $\pi$. Then
\[
\hat T^{(K)}=\tW(\tW^\top\tW)^r+\sum_{\pi:\,\exists k_e>1}\ \sum_{\sigma\ge\pi}\mu(\pi,\sigma)\big(g^{\mathrm{corr}}_\pi(\sigma)-g^{\mathrm{plug}}_\pi(\sigma)\big),
\]
with $\mu$ the Möbius function of the product of the two partition lattices,
satisfies $\E\hat T^{(K)}=(\tY\tY^\top)^r\tY$ under the conditions of
Theorem~\ref{thm:unbiased} with $\xi\equiv0$.
\end{theorem}
The number of pairs $(\pi,\sigma)$ is 5, 99 and 2{,}790 for $K=3,5,7$, and for
$K=3$ the construction reduces to Eq.~\eqref{eq:drup}. Theorem~\ref{thm:unbiased},
Corollary~\ref{cor:cand} and Theorem~\ref{thm:fair} carry over to any odd $K$.

Assumption~\ref{as:unconf} excludes slates of fixed size, where exposures in a
user's history are negatively dependent (Coat, for instance, collects exactly
24 ratings per user). The next result bounds the resulting bias.

\begin{proposition}[Exposure dependence within users]\label{prop:dep}
Suppose that exposures are independent across users but may be dependent within
a user's row, with $|\mathrm{Cov}(O_e,O_f\mid Y)|\le\rho\,p_ep_f$ for distinct $e,f$ in
the same row. With correct propensities,
\[
\big|\E\hat T_{ui}-T^*_{ui}\big|\ \le\ \rho\sum_{j\neq i}C_{uj}\Big(\sum_{v\neq u}C_{vj}C_{vi}\,\varepsilon_{vj}\varepsilon_{vi}+C_{uj}C_{ui}\,\varepsilon_{uj}\varepsilon_{ui}\Big).
\]
For slates of $k$ items drawn uniformly without replacement, $\rho\le1/k$, and
$\rho\le1$ for any negatively dependent design. For DR edges the bias is of
second order in the imputation error; for IPS edges ($\varepsilon=Y$) it is of
first order.
\end{proposition}

\begin{theorem}[Repeated-walk bias of uncorrected propagation]\label{thm:lower}
Let the propensities be correct and let $(u,i)$ be a candidate with $O_{ui}=0$.
Relative to the conditional target of Corollary~\ref{cor:cand}, the
uncorrected DR operator $\tW\tW^\top\tW$ has conditional bias
\begin{gather*}
C_{ui}\hat Y_{ui}\big(R_u^{(i)}+K_i\big)-C_{ui}^3\hat Y_{ui}(1-\hat Y_{ui}^2),\\
K_i=\sum_{v\neq u}C_{vi}^2V_{vi},\qquad R_u^{(i)}=\sum_{j\neq i}C_{uj}^2V_{uj},
\end{gather*}
where $V_e=\Var(W_e)=\frac{1-p_e}{p_e}\varepsilon_e^2$. Moreover:
(a) $K_i\ge(1/\max_vp_{vi}-1)\sum_{v\neq u}C_{vi}^2\varepsilon_{vi}^2$, which is unbounded
as the item's propensities vanish;
(b) with clipped propensities $\bar p=\max(p,\tau)$, $\tau\le1/2$, the excess
$C_e^2(\E[W_e^2]-\E[W_e])$ of a doubly traversed edge is at most
$C_e^2(1/\tau-1)\varepsilon_e^2$, with equality at $p_e=\tau$; the worst-case
repeated-walk bias is therefore $\Theta(1/\tau)$;
(c) for IPS edges the conditional bias vanishes at three hops.
DRUP has zero bias relative to the same target.
\end{theorem}
The within-user ranking is affected through $C_{ui}\hat Y_{ui}K_i$, which
up-weights items whose other raters were rarely exposed; $R_u^{(i)}$ depends on
$i$ only through one term.

\begin{theorem}[Variance, clipping and concentration]\label{thm:var}
Let $\Lambda_e(u,i)$ be the $C$-weighted mass of the walks from $u$ to $i$ through
edge $e$, and $B=\max(1,1/\tau)$. Then
(a) $\Var(W_e)\le\varepsilon_e^2/(4\tau^2)$, and $\Var(W_e)\le\varepsilon_e^2/\tau$ if $\hat p=p$;
(b) if $\hat p=p$, $\Var(\hat T_{ui})\le\frac{(1+1/\tau)^2}{\tau}\sum_e\varepsilon_e^2\Lambda_e(u,i)^2$;
(c) if $\hat p=p$, so that clipping is the only source of bias,
$|\E\hat T_{ui}-T^*_{ui}|\le3\bar\xi(1+\bar\xi)^2\sum_{j,v}C_{uj}C_{vj}C_{vi}$ with
$\bar\xi=\max_e(1-p_e/\tau)_+\varepsilon_e$;
(d) for any $\hat p$, $\Pr(|\hat T_{ui}-\E\hat T_{ui}|\ge t)\le2\exp(-2t^2/\sum_ec_e^2)$ with
$c_e=\varepsilon_eB^2\Lambda_e(u,i)/\tau$.
\end{theorem}

The variance scales with the imputation error $\varepsilon^2$, not with $Y^2$ as
for IPS, and $\tau$ trades variance against a clipping bias that vanishes when
$\hat Y$ is accurate. The constants are worst-case and loose; Section~\ref{sec:mc}
compares bound (b) with the Monte-Carlo variance.

\begin{definition}[Exposure elasticity]\label{def:elast}
For an intervention $do(p_{\cdot i}\leftarrow\pi p_{\cdot i})$ on item $i$'s
exposure mechanism, after which the estimator uses the post-intervention
propensities ($\hat p_{\cdot i}\leftarrow\pi\hat p_{\cdot i}$) and otherwise the same
nuisances, $\eta_i=\partial\log\E[s_{ui}]/\partial\log\pi$.
\end{definition}

\begin{theorem}[Exposure invariance]\label{thm:fair}
Under Assumptions~\ref{as:unconf}--\ref{as:fixed} and $\pi p\ge\tau$, DRUP has
$\eta_i=0$, also conditionally on $O_{ui}=0$: the expected score of every
candidate is unchanged when item $i$'s exposure mechanism is scaled.
Logged-graph propagation with fixed weights $C$ has $\eta_i\ge1$ (popularity
amplification); normalising by logged degrees partly offsets it.
Uncorrected DR propagation has $\eta_i<0$ whenever $\hat Y_{ui}K_i>0$
(over-correction).
\end{theorem}

Exposure invariance is a statement about expected scores. It is not a
distributive fairness criterion, and it is conditional: when many propensities
lie below the clip, $\pi p\ge\tau$ fails and invariance is lost
(Section~\ref{sec:fair}).

\subsection{Structural properties}\label{sec:structsum}

Given its nuisances, $s$ is a fixed linear function of the edge estimates.
Section~\ref{sec:structth} of the Supplementary Material states and proves the
consequences, which hold for every linear propagation operator, corrected or
not: exact attributions and minimal counterfactual explanations with frozen
nuisances (Proposition~\ref{prop:affine}, Theorem~\ref{thm:cf}), conservative
certificates against injected profiles (Theorem~\ref{thm:robust}), joint
differential privacy when the population-level nuisances, the
hyper-parameters and the score constants are public (Theorem~\ref{thm:dp}),
and an exposure-capped allocation solved exactly by min-cost flow
(Theorem~\ref{thm:alloc}). Because they are not specific to DRUP, we report
their empirical behaviour in the Supplementary Material
(Section~\ref{sec:struct}) and use only the allocation in the main text, as a
decision that compares scores across users (Section~\ref{sec:semisynth}).
